\chapter{Architektur und Entwurfsentscheidungen}
\label{chap:architektur-entwurf}

\section{Gesamtarchitektur und Schichtenmodell}
\label{sec:systemarchitektur-ueberblick}
Die Architektur der Gesamtlösung wurde nach den Richtlinien des AWS Well-Architected Frameworks (Reliability Pillar) entworfen und gliedert sich in drei klar voneinander abgegrenzte funktionale Schichten:\customcite{aws_reliability_2025}
\begin{enumerate}[noitemsep,topsep=2pt]
    \item \textbf{Infrastrukturschicht (AWS):} Bildet das cloud-basierte Netzwerk- und Rechenfundament. Eine dedizierte Virtual Private Cloud (VPC, CIDR \texttt{10.0.0.0/16}) kapselt die Umgebung ab. Ein öffentliches Subnetz (\texttt{10.0.1.0/24}) in der Region Frankfurt (\texttt{eu-central-1}) ist über ein Internet Gateway angebunden. Eine EC2-Instanz des Typs \texttt{t3.small} (2 vCPUs, 2~GB RAM, 20~GB gp3-SSD) stellt die Rechenleistung bereit. Eine statische Elastic IP sichert die persistente Erreichbarkeit über Instanzneustarts hinweg.
    \item \textbf{Plattformschicht (k3s \& Argo~CD):} k3s stellt die CNCF-zertifizierte Kubernetes-Laufzeitumgebung bereit. Alle Control-Plane- und Worker-Komponenten laufen ressourceneffizient auf dem Knoten. Als GitOps-Controller ist Argo~CD im Namespace \texttt{argocd} installiert und überwacht das GitHub-Repository.
    \item \textbf{An\-wen\-dungs\-schicht (Resilience Lab):} Beherbergt die modularisierte Demo-Webanwendung, bestehend aus der zustandslosen REST-API, einem dynamischen Web-Dashboard und geschützten Endpunkten zur Fehlerinduktion.
\end{enumerate}

Das Zusammenspiel dieser drei Architekturschichten, die Netzwerkgrenzen sowie die Signal- und GitOps-Steuerungsflüsse sind in Abbildung~\ref{fig:architektur-uebersicht} visualisiert.

\begin{figure}[htbp]
\centering
\begin{tikzpicture}[
    font=\sffamily\footnotesize,
    node distance=0.8cm and 1.0cm,
    ext/.style={draw=gray!70, fill=gray!12, rounded corners=3pt, minimum width=2.7cm, minimum height=0.85cm, align=center, font=\sffamily\scriptsize},
    k8scomp/.style={draw=blue!70!black, fill=blue!8, rounded corners=3pt, minimum width=2.4cm, minimum height=0.8cm, align=center, font=\sffamily\scriptsize},
    pod/.style={draw=teal!70!black, fill=teal!10, rounded corners=3pt, minimum width=2.8cm, minimum height=0.9cm, align=center, font=\sffamily\scriptsize},
    arrow/.style={->, >=stealth, line width=0.8pt, draw=black!80},
    gitflow/.style={->, >=stealth, line width=0.8pt, dashed, draw=blue!70!black}
]

% 1. Externe Komponenten
\node[ext] (client) {\textbf{Client / Lastgenerator}\\HTTP Traffic (15 Req/s)};
\node[ext, right=1.6cm of client] (github) {\textbf{GitHub Repository}\\\path{k3s_aws (GitOps)}\\{\tiny Source of Truth (Overlays)}};
\node[ext, right=1.6cm of github] (tf) {\textbf{HashiCorp Terraform}\\{\tiny Day-0 IaC Provisioning}};

% 2. Cluster interne Komponenten
\node[k8scomp, below=2.4cm of client] (traefik) {\textbf{Traefik Ingress}\\{\tiny Port 80 / 443}};
\node[k8scomp, right=1.0cm of traefik] (service) {\textbf{ClusterIP Service}\\\texttt{resilience-lab:8000}};
\node[k8scomp, right=1.0cm of service] (argocd) {\textbf{Argo CD Controller}\\{\tiny Sync \& Self-Heal Loop}};

% Pods
\node[pod, below=0.9cm of service, xshift=-1.6cm] (pod1) {\textbf{Pod 1: \texttt{resilience-lab}}\\{\tiny FastAPI App (\texttt{1/1 Ready})}\\{\tiny Health: \texttt{/health/live}, \texttt{/ready}}};
\node[pod, below=0.9cm of service, xshift=1.6cm] (pod2) {\textbf{Pod 2: \texttt{resilience-lab}}\\{\tiny FastAPI App (\texttt{1/1 Ready})}\\{\tiny Health: \texttt{/health/live}, \texttt{/ready}}};

% 3. K3s Box
\node[draw=teal!60!black, fill=teal!3, rounded corners=4pt, dashed, line width=0.7pt,
      fit=(traefik) (argocd) (pod1) (pod2),
      inner sep=10pt,
      label={[anchor=north west, font=\sffamily\bfseries\scriptsize, text=teal!70!black]north west:\textbf{k3s Kubernetes Cluster} (Control Plane \& Worker in einer Binärdatei, SQLite WAL)}] (k3sbox) {};

% 4. EC2 Box
\node[draw=orange!70!black, fill=orange!3, rounded corners=5pt, line width=0.8pt,
      fit=(k3sbox),
      inner xsep=10pt, inner ysep=14pt,
      label={[anchor=north west, font=\sffamily\bfseries\scriptsize, text=orange!80!black]north west:\textbf{AWS EC2 Instanz} (\texttt{t3.small}, Elastic IP: \texttt{63.176.27.134}, Ubuntu 24.04 LTS)}] (ec2box) {};

% 5. Subnetz Box
\node[draw=gray!50, fill=white, rounded corners=6pt, line width=0.6pt,
      fit=(ec2box),
      inner xsep=8pt, inner ysep=14pt,
      label={[anchor=north west, font=\sffamily\bfseries\scriptsize, text=gray!70!black]north west:Öffentliches Subnetz (\texttt{10.0.1.0/24}) \textbar{} Security Group (Ports: 22, 80, 443, 6443)}] (subnetbox) {};

% 6. VPC / AWS Cloud Box
\node[draw=blue!60!black, fill=blue!2, rounded corners=7pt, line width=0.8pt,
      fit=(subnetbox),
      inner xsep=8pt, inner ysep=15pt,
      label={[anchor=north west, font=\sffamily\bfseries\scriptsize, text=blue!80!black]north west:\textbf{AWS Cloud (Region: \texttt{eu-central-1} Frankfurt)} \textbar{} VPC: \texttt{10.0.0.0/16} \textbar{} Internet Gateway}] (awsbox) {};

% Pfeile und Verbindungen
\draw[arrow] (client) -- node[right, font=\tiny] {HTTP:80} (client |- awsbox.north);
\draw[arrow] (client |- awsbox.north) -- (traefik);
\draw[arrow] (traefik) -- node[above, font=\tiny] {Reverse Proxy} (service);
\draw[arrow] (service.south) -- ++(0,-0.3) -| (pod1.north);
\draw[arrow] (service.south) -- ++(0,-0.3) -| (pod2.north);

\draw[gitflow] (github) -- node[right, font=\tiny] {Git Pull Manifests} (github |- awsbox.north);
\draw[gitflow] (github |- awsbox.north) -- (argocd);
\draw[gitflow] (argocd.south) -- ++(0,-0.3) -| node[pos=0.3, below, font=\tiny] {Reconcile \& Self-Heal} (pod2.north);

\draw[arrow, dashed, draw=orange!80!black] (tf) -- node[right, font=\tiny] {Day-0 Cloud-Init} (tf |- awsbox.north);

\end{tikzpicture}
\caption{System- und Netzwerkarchitektur der GitOps-gesteuerten Kubernetes-Umgebung auf AWS, eigene Darstellung.}
\label{fig:architektur-uebersicht}
\end{figure}


\section{Wirtschaftliche und technische Bewertung: k3s gegenüber Amazon EKS}
\label{sec:wahl-k3s-vs-eks}
Ein zentraler architektonischer Entwurfsentscheid betrifft die Auswahl der Kubernetes-Plattform. Enterprise-Umgebungen greifen häufig standardmäßig auf Amazon Elastic Kubernetes Service (EKS) zurück. Für Entwicklungs-, Staging- oder universitäre Forschungsumgebungen weist EKS jedoch ein ungünstiges Kosten-Nutzen-Verhältnis auf:
\begin{itemize}[noitemsep,topsep=2pt]
    \item \textbf{Amazon EKS:} Allein für die Bereitstellung des EKS-Cluster-Endpunkts fallen feste Basiskosten von 0{,}10~USD pro Stunde (ca. 73~USD/Monat) an. Für Worker-Nodes (mindestens zwei \texttt{t3.medium}-Instanzen für Hochverfügbarkeit der Control Plane), einen AWS Application Load Balancer (ca. 22~USD/Monat) und NAT-Gateways belaufen sich die monatlichen Gesamtkosten auf über 180 bis 220~USD.
    \item \textbf{k3s auf EC2 (\textnormal{t3.small}):} Die Nutzung einer einzelnen reservierten oder On-Demand-\texttt{t3.small}-Instanz verursacht lediglich ca. 0{,}026~USD pro Stunde (ca. 19~USD/Monat). Der integrierte Traefik-Ingress-Controller erfordert keinen teuren AWS-Load-Balancer, sondern bindet direkt an die Ports 80 und 443 der Elastic IP.
\end{itemize}

Da die in dieser Arbeit zu evaluierenden Mechanismen (Prozessneustarts, Readiness-Gating und GitOps-Self-Heal) auf Pod-, Service- und Ingress-Ebene innerhalb der standardisierten Kubernetes-API operieren, liefert das k3s-Setup identische empirische Erkenntnisse bei einer Kostenreduktion von ca. 90\,\%.

\section{Zuständigkeitsabgrenzung: Terraform und Argo~CD}
\label{sec:zustaendigkeit-terraform-argocd}
Um konkurrierende Schreibzugriffe und Race Conditions zwischen verschiedenen Automatisierungswerkzeugen zu verhindern, folgt die Architektur einem strikten Separationsprinzip:
\begin{itemize}[noitemsep,topsep=2pt]
    \item \textbf{Terraform (Day-0-Infrastruktur):} Verantwortlich ausschließlich für das Provisionieren der Cloud-Ressourcen (VPC, Subnetze, Routing, Security Groups, EC2, Elastic IP) sowie für das einmalige Bootstrappen des Betriebssystems via Cloud-Init. Nach erfolgreichem Aufsetzen von k3s und Argo~CD verbleibt der Terraform-State unverändert.
    \item \textbf{Argo~CD (Day-2-Anwendungsbetrieb):} Übernimmt ab dem Bootstrap die alleinige Verantwortung für alle Kubernetes-Ressourcen im Cluster. Weder Entwickler noch CI-Systeme greifen mit imperativen \texttt{kubectl}-Befehlen modifizierend in die Produktivumgebung ein.
\end{itemize}

\section{Manifest-Strukturierung mit Kustomize und Workflow-Trennung}
\label{sec:kustomize-manifestverwaltung}
Für die Konfigurationsverwaltung wird Kustomize eingesetzt. Im Gegensatz zu Helm, das komplexe Go-Templates und Tiller-artige Paketierungen verwendet, verfolgt Kustomize einen rein deklarativen Ansatz (Plain YAML). Es nutzt ein hierarchisches Base-Overlay-Muster:
\begin{itemize}[noitemsep,topsep=2pt]
    \item \texttt{kustomize/base}: Definiert den kanonischen Minimalzustand der Anwendung (Deployment mit Standard-Labels, Container-Port 8000, Liveness- und Readiness-Probes sowie ein ClusterIP-Service).
    \item \texttt{kustomize/overlays/local}: Patcht die Konfiguration für lokale Entwicklungscluster (z.\,B. NodePort 30080 und lokales Image-Pulling).
    \item \texttt{kustomize/overlays/aws}: Erhöht die Replika-Zahl auf 2 für Resilienzprüfungen, setzt Produktions-Image-Tags aus der GitHub Container Registry und fügt ein Ingress-Objekt für den Traefik-Controller hinzu.
\end{itemize}

Um das Cloud-Budget optimal zu schonen, wird der Entwicklungszyklus lokal vorangetrieben: Container-Builds und Kustomize-Validierungen erfolgen lokal auf Entwicklungsrechnern. Die AWS-Infrastruktur wird gezielt und automatisiert über Terraform erst dann hochgefahren, wenn die Experimente unter realen Netzwerkbedingungen evaluiert werden.

\section{Sicherheits- und Netzwerkkonzept}
\label{sec:sicherheitskonzept}
Die Netzwerksicherheit stützt sich auf eine strikte Trennung von externen und internen Verkehrsflüssen über AWS Security Groups:
\begin{itemize}[noitemsep,topsep=2pt]
    \item \textbf{Administrativer Zugang (SSH \& K8s-API):} Der SSH-Zugriff über Port 22 ist ausschließlich über asymmetrische RSA-Schlüsselpaare (2048 Bit) möglich; Passwort-Authentifizierung ist serverseitig deaktiviert. Die Kubernetes-API (Port 6443) wird über mTLS-Zertifikate abgesichert, die während des Bootstraps mit der Elastic IP als Subject Alternative Name (SAN) generiert werden.
    \item \textbf{Anwendungs-Ingress (HTTP/HTTPS):} Der externe Produktivverkehr wird ausschließlich über die Ports 80 und 443 an den Traefik-Ingress-Controller weitergeleitet. Sämtliche containerinternen Ports (wie der FastAPI-Port 8000) bleiben vollständig im virtuellen Overlay-Netzwerk (Flannel) des Clusters isoliert und sind von außen zu keinem Zeitpunkt direkt adressierbar.
\end{itemize}
